\documentclass[11pt,a4paper,notitlepage]{report}
% Shared preamble for the raw chapter-4 write-ups under docs/report/insights/.
%
% Each experiment_N/experiment_N.tex is a standalone, compilable LaTeX report
% written BEFORE anything goes into final_report/chapters/chapter_4_*.tex.
% It uses the final report's fonts, units and macros so that passages, tables
% and figures lift into the chapter without reformatting.
%
%   \documentclass[11pt,a4paper]{report}
%   % Shared preamble for the raw chapter-4 write-ups under docs/report/insights/.
%
% Each experiment_N/experiment_N.tex is a standalone, compilable LaTeX report
% written BEFORE anything goes into final_report/chapters/chapter_4_*.tex.
% It uses the final report's fonts, units and macros so that passages, tables
% and figures lift into the chapter without reformatting.
%
%   \documentclass[11pt,a4paper]{report}
%   % Shared preamble for the raw chapter-4 write-ups under docs/report/insights/.
%
% Each experiment_N/experiment_N.tex is a standalone, compilable LaTeX report
% written BEFORE anything goes into final_report/chapters/chapter_4_*.tex.
% It uses the final report's fonts, units and macros so that passages, tables
% and figures lift into the chapter without reformatting.
%
%   \documentclass[11pt,a4paper]{report}
%   \input{../common/preamble}
%   \begin{document} ...
%
% Figures and tables are NOT copied: \graphicspath and \reporttable resolve into
% final_report/figures/ and final_report/tables/, so the raw write-up always
% shows the same artefacts as the chapter.  Compile with:  latexmk -pdf <file>

\usepackage[T1]{fontenc}
\usepackage[utf8]{inputenc}
\usepackage{mathptmx}
\usepackage[scaled=0.92]{helvet}
\usepackage{courier}
\usepackage[a4paper,left=30mm,right=20mm,top=25mm,bottom=25mm]{geometry}
\usepackage{microtype}
\usepackage{booktabs}
\usepackage{tabularx}
\usepackage{longtable}
\usepackage{array}
\usepackage{multirow}
\usepackage{amsmath}
\usepackage{amssymb}
\usepackage{ifthen}
\usepackage{siunitx}
\usepackage{graphicx}
\usepackage{placeins}
\usepackage{caption}
\usepackage{enumitem}
\usepackage[table]{xcolor}
\usepackage{rotating}
\PassOptionsToPackage{hyphens}{url}\usepackage{url}
\usepackage[hidelinks]{hyperref}
\usepackage[nameinlink,noabbrev]{cleveref}

\graphicspath{{../../final_report/}{../../final_report/figures/}}
% \reporttable{name} inputs final_report/tables/name.tex
\newcommand{\reporttable}[1]{\input{../../final_report/tables/#1}}

\newcommand{\repo}[1]{\texttt{\detokenize{#1}}}
\sisetup{detect-all,group-separator={,},range-phrase=--,range-units=single}
\DeclareSIUnit{\pixel}{px}
\DeclareSIUnit{\mebi}{Mi}
\captionsetup{font=small,labelfont=bf}
\setlist{nosep}
\setlength{\parskip}{4pt}
\setlength{\LTcapwidth}{\textwidth}

% ---- result-slot status (RENDERED here, unlike in the final report)
\definecolor{stcompleted}{HTML}{A5D6A7}
\definecolor{stpartial}{HTML}{FFF2CC}
\definecolor{stuntested}{HTML}{FFD8A8}
\definecolor{stblocked}{HTML}{FFCDD2}
\definecolor{stfailed}{HTML}{FFCDD2}
\definecolor{stplanned}{HTML}{E0E0E0}
\newcommand{\slotstatus}[2]{%
  \par\noindent\fcolorbox{black!40}{st#1}{\parbox{\dimexpr\textwidth-2\fboxsep-2\fboxrule}{%
    \textbf{Status: #1.}\ #2}}\par\medskip}
\newcommand{\status}[1]{\textbf{#1}}

% ---- queries: RENDERED in red here so a reader sees every open concern in place
\definecolor{flagred}{HTML}{B00000}
\newcommand{\flag}[1]{\textcolor{flagred}{#1}}
\newcommand{\moved}[1]{#1}
\newenvironment{queried}[1][Questionable]%
  {\par\addvspace{0.6ex}\begingroup\color{flagred}\noindent
   \textbf{[#1]}\ \ignorespaces}%
  {\par\endgroup\addvspace{0.6ex}}

% ---- per-block performance marks (section 4.5 diagrams)
\definecolor{perfgood}{HTML}{A5D6A7}
\definecolor{perfpoor}{HTML}{FFCDD2}
\definecolor{perfnoeffect}{HTML}{D9EAF7}
\definecolor{perfnone}{HTML}{FFF2CC}
\newcommand{\perfmark}[1]{%
  \ifthenelse{\equal{#1}{good}}{\colorbox{perfgood}{\scriptsize good}}{%
  \ifthenelse{\equal{#1}{poor}}{\colorbox{perfpoor}{\scriptsize poor}}{%
  \ifthenelse{\equal{#1}{noeffect}}{\colorbox{perfnoeffect}{\scriptsize no effect}}{%
  \colorbox{perfnone}{\scriptsize no claim}}}}}

%   \begin{document} ...
%
% Figures and tables are NOT copied: \graphicspath and \reporttable resolve into
% final_report/figures/ and final_report/tables/, so the raw write-up always
% shows the same artefacts as the chapter.  Compile with:  latexmk -pdf <file>

\usepackage[T1]{fontenc}
\usepackage[utf8]{inputenc}
\usepackage{mathptmx}
\usepackage[scaled=0.92]{helvet}
\usepackage{courier}
\usepackage[a4paper,left=30mm,right=20mm,top=25mm,bottom=25mm]{geometry}
\usepackage{microtype}
\usepackage{booktabs}
\usepackage{tabularx}
\usepackage{longtable}
\usepackage{array}
\usepackage{multirow}
\usepackage{amsmath}
\usepackage{amssymb}
\usepackage{ifthen}
\usepackage{siunitx}
\usepackage{graphicx}
\usepackage{placeins}
\usepackage{caption}
\usepackage{enumitem}
\usepackage[table]{xcolor}
\usepackage{rotating}
\PassOptionsToPackage{hyphens}{url}\usepackage{url}
\usepackage[hidelinks]{hyperref}
\usepackage[nameinlink,noabbrev]{cleveref}

\graphicspath{{../../final_report/}{../../final_report/figures/}}
% \reporttable{name} inputs final_report/tables/name.tex
\newcommand{\reporttable}[1]{\input{../../final_report/tables/#1}}

\newcommand{\repo}[1]{\texttt{\detokenize{#1}}}
\sisetup{detect-all,group-separator={,},range-phrase=--,range-units=single}
\DeclareSIUnit{\pixel}{px}
\DeclareSIUnit{\mebi}{Mi}
\captionsetup{font=small,labelfont=bf}
\setlist{nosep}
\setlength{\parskip}{4pt}
\setlength{\LTcapwidth}{\textwidth}

% ---- result-slot status (RENDERED here, unlike in the final report)
\definecolor{stcompleted}{HTML}{A5D6A7}
\definecolor{stpartial}{HTML}{FFF2CC}
\definecolor{stuntested}{HTML}{FFD8A8}
\definecolor{stblocked}{HTML}{FFCDD2}
\definecolor{stfailed}{HTML}{FFCDD2}
\definecolor{stplanned}{HTML}{E0E0E0}
\newcommand{\slotstatus}[2]{%
  \par\noindent\fcolorbox{black!40}{st#1}{\parbox{\dimexpr\textwidth-2\fboxsep-2\fboxrule}{%
    \textbf{Status: #1.}\ #2}}\par\medskip}
\newcommand{\status}[1]{\textbf{#1}}

% ---- queries: RENDERED in red here so a reader sees every open concern in place
\definecolor{flagred}{HTML}{B00000}
\newcommand{\flag}[1]{\textcolor{flagred}{#1}}
\newcommand{\moved}[1]{#1}
\newenvironment{queried}[1][Questionable]%
  {\par\addvspace{0.6ex}\begingroup\color{flagred}\noindent
   \textbf{[#1]}\ \ignorespaces}%
  {\par\endgroup\addvspace{0.6ex}}

% ---- per-block performance marks (section 4.5 diagrams)
\definecolor{perfgood}{HTML}{A5D6A7}
\definecolor{perfpoor}{HTML}{FFCDD2}
\definecolor{perfnoeffect}{HTML}{D9EAF7}
\definecolor{perfnone}{HTML}{FFF2CC}
\newcommand{\perfmark}[1]{%
  \ifthenelse{\equal{#1}{good}}{\colorbox{perfgood}{\scriptsize good}}{%
  \ifthenelse{\equal{#1}{poor}}{\colorbox{perfpoor}{\scriptsize poor}}{%
  \ifthenelse{\equal{#1}{noeffect}}{\colorbox{perfnoeffect}{\scriptsize no effect}}{%
  \colorbox{perfnone}{\scriptsize no claim}}}}}

%   \begin{document} ...
%
% Figures and tables are NOT copied: \graphicspath and \reporttable resolve into
% final_report/figures/ and final_report/tables/, so the raw write-up always
% shows the same artefacts as the chapter.  Compile with:  latexmk -pdf <file>

\usepackage[T1]{fontenc}
\usepackage[utf8]{inputenc}
\usepackage{mathptmx}
\usepackage[scaled=0.92]{helvet}
\usepackage{courier}
\usepackage[a4paper,left=30mm,right=20mm,top=25mm,bottom=25mm]{geometry}
\usepackage{microtype}
\usepackage{booktabs}
\usepackage{tabularx}
\usepackage{longtable}
\usepackage{array}
\usepackage{multirow}
\usepackage{amsmath}
\usepackage{amssymb}
\usepackage{ifthen}
\usepackage{siunitx}
\usepackage{graphicx}
\usepackage{placeins}
\usepackage{caption}
\usepackage{enumitem}
\usepackage[table]{xcolor}
\usepackage{rotating}
\PassOptionsToPackage{hyphens}{url}\usepackage{url}
\usepackage[hidelinks]{hyperref}
\usepackage[nameinlink,noabbrev]{cleveref}

\graphicspath{{../../final_report/}{../../final_report/figures/}}
% \reporttable{name} inputs final_report/tables/name.tex
\newcommand{\reporttable}[1]{\input{../../final_report/tables/#1}}

\newcommand{\repo}[1]{\texttt{\detokenize{#1}}}
\sisetup{detect-all,group-separator={,},range-phrase=--,range-units=single}
\DeclareSIUnit{\pixel}{px}
\DeclareSIUnit{\mebi}{Mi}
\captionsetup{font=small,labelfont=bf}
\setlist{nosep}
\setlength{\parskip}{4pt}
\setlength{\LTcapwidth}{\textwidth}

% ---- result-slot status (RENDERED here, unlike in the final report)
\definecolor{stcompleted}{HTML}{A5D6A7}
\definecolor{stpartial}{HTML}{FFF2CC}
\definecolor{stuntested}{HTML}{FFD8A8}
\definecolor{stblocked}{HTML}{FFCDD2}
\definecolor{stfailed}{HTML}{FFCDD2}
\definecolor{stplanned}{HTML}{E0E0E0}
\newcommand{\slotstatus}[2]{%
  \par\noindent\fcolorbox{black!40}{st#1}{\parbox{\dimexpr\textwidth-2\fboxsep-2\fboxrule}{%
    \textbf{Status: #1.}\ #2}}\par\medskip}
\newcommand{\status}[1]{\textbf{#1}}

% ---- queries: RENDERED in red here so a reader sees every open concern in place
\definecolor{flagred}{HTML}{B00000}
\newcommand{\flag}[1]{\textcolor{flagred}{#1}}
\newcommand{\moved}[1]{#1}
\newenvironment{queried}[1][Questionable]%
  {\par\addvspace{0.6ex}\begingroup\color{flagred}\noindent
   \textbf{[#1]}\ \ignorespaces}%
  {\par\endgroup\addvspace{0.6ex}}

% ---- per-block performance marks (section 4.5 diagrams)
\definecolor{perfgood}{HTML}{A5D6A7}
\definecolor{perfpoor}{HTML}{FFCDD2}
\definecolor{perfnoeffect}{HTML}{D9EAF7}
\definecolor{perfnone}{HTML}{FFF2CC}
\newcommand{\perfmark}[1]{%
  \ifthenelse{\equal{#1}{good}}{\colorbox{perfgood}{\scriptsize good}}{%
  \ifthenelse{\equal{#1}{poor}}{\colorbox{perfpoor}{\scriptsize poor}}{%
  \ifthenelse{\equal{#1}{noeffect}}{\colorbox{perfnoeffect}{\scriptsize no effect}}{%
  \colorbox{perfnone}{\scriptsize no claim}}}}}

\setcounter{chapter}{4}
\renewcommand{\thesection}{\arabic{section}}
\renewcommand{\thefigure}{I.\arabic{figure}}
\renewcommand{\thetable}{I.\arabic{table}}
% every table and figure of this document is generated by script/exp1_insights.py
\newcommand{\insdir}{../../../../output/compare/simulation/exp1-insights}
\newcommand{\ins}[1]{\input{\insdir/#1}}

\title{Experiment 1: Why the Estimators Err the Way They Do\\
       \large Insight analysis for Chapter 4, Section 4.1}
\author{Bolton Athit Davies}
\date{Reporting cutoff: 22 October 2026}

\begin{document}
\maketitle

\noindent\textbf{Purpose of this document.} The raw write-up of Experiment~1
(\path{insights/experiment_1/}) reports how far each estimator is from the
truth on the five static recordings. This document asks what kind of error
each one makes and when it is made: whether wheel odometry's error is a scale
or a heading error, why IMU dead reckoning diverges, when and why the EKF's
covariance stops describing its error, and what the visual estimators' error
consists of once they are scored at the IMU point. Each question is tested
within every recording, so the recording is its own control, and the answer
carries the evidence for and against it. Every number is written by
\path{script/exp1_insights.py} into
\path{output/compare/simulation/exp1-insights/}; the tables are input from
there unchanged.

\section{Questions and method}

Six questions were put to the five \repo{dataset_allsensor} recordings
(replay at $0.5\times$, four worlds, one run per estimator per recording).

\begin{enumerate}[label=Q\arabic*.]
\item Is wheel odometry's error a distance (scale) error, a heading error, or
      heading error carried into position?
\item Does each estimator's error grow faster while the robot turns than on
      straights?
\item Does IMU dead reckoning diverge through its heading or through its
      velocity?
\item When does the EKF's covariance stop describing its error, and does the
      per-recording noise calibration change that?
\item What does the transferred noise model (the VSLAM file's IMU constants)
      do to the EKF, mechanically?
\item For the visual estimators, scored at the IMU point: what is their error
      made of, does it grow at turns or with few features, and what was the one
      failure?
\end{enumerate}

\paragraph{Baselines.} The wheel, IMU and EKF estimates start from the
ground-truth pose and are compared to the raw model-origin ground truth, as in
the raw write-up. The error vector is rotated into the ground-truth path
frame, giving an along-track part (distance travelled) and a cross-track part
(lateral), and the heading error is the estimate's yaw against the truth's,
zeroed at the start. To separate heading from velocity error, the true
velocity is integrated through the heading error alone; the ratio of that
prediction to the actual error says how much of the error the heading
explains. Growth rates while turning (ground-truth yaw rate above
\SI{10}{\degree\per\second}) and on straights are compared as in the
Experiment~3 insight analysis. NEES is computed from the logged covariance at
\SI{1}{\hertz}; the filter ``leaves its band'' at the first time it stays
above the \SI{97.5}{\percent} value of a 3-DOF $\chi^2$ for \SI{5}{\second}.

\paragraph{Visual estimators.} VINS-Fusion and ORB-SLAM3 report the IMU,
\SI{0.338}{\metre} ahead of the model origin, so they are scored against the
ground truth moved to that point (the lever-arm finding of the Experiment~3
insight analysis, now applied in every scoring script). The initialisation
transient is cut and the start-anchored error, tracking regime and
turning contrast of the Experiment~3 analysis are reused. Whole-run ATE is
reported both in three dimensions, as the chapter does, and in the plane.

\FloatBarrier
\section{Insight 1: wheel odometry's error is heading error made in the
         turns, not a distance error}
\label{sec:insight1-wheel}

Wheel odometry measures distance to within one percent; its position error is
heading error carried forward, and the heading error is made almost entirely
while the robot turns. \Cref{tab:e1i-decomp,tab:e1i-turning} and
\cref{fig:decomp} carry the evidence.

\ins{decomposition.tex}
\ins{turning.tex}

\begin{figure}[p]
  \centering
  \includegraphics[width=\textwidth]{\insdir/fig_decomposition_light.png}
  \caption{Anchored error of wheel odometry and the EKF on each recording, in
  the ground-truth path frame: along-track error, cross-track error and heading
  error against time. Grey bands are the moments the ground-truth yaw rate
  exceeds \SI{10}{\degree\per\second}.}
  \label{fig:decomp}
\end{figure}

The path length the wheels report is within \SI{0.8}{\percent} of the true one
on every recording ($n=5$), while the heading error reaches
\SIrange{7}{16}{\degree} RMS. Integrating the true velocity through that
heading error alone reproduces \SIrange{95}{104}{\percent} of the final
position error on all five recordings. The along-track component is not a
scale error: it appears once the heading error has rotated the path, which is
why it is as large as the cross-track component.

The heading error is made in the turns. Its growth rate while turning is
\SIrange{650}{1406}{\milli\degree\per\second} against a negative or near-zero
rate on straights, in five of five recordings; the heading-error traces of
\cref{fig:decomp} step at each grey band and hold between them. The size of
the step varies between turns and between recordings, which is why the raw
write-up found no relation between the fitted yaw-rate scale and the total
drift (its H1); the mechanism is nonetheless a per-turn heading error, not a
continuous drift.

\FloatBarrier
\section{Insight 2: IMU dead reckoning diverges through its velocity, not
         its heading}

IMU-only dead reckoning keeps its heading to within a few degrees and loses
its position through the velocity integral, whose error grows roughly as
$t^{1.5}$. \Cref{tab:e1i-imu} and \cref{fig:imu} carry the evidence.

\ins{imu_law.tex}

\begin{figure}[htbp]
  \centering
  \includegraphics[width=\textwidth]{\insdir/fig_imu_law_light.png}
  \caption{IMU dead reckoning on the five recordings. Left: anchored position
  error against time on logarithmic axes, with $t$ and $t^{2}$ guides. Right:
  the same error on linear axes with, dashed, the error that the heading error
  alone would produce.}
  \label{fig:imu}
\end{figure}

The heading error of the IMU integration is \SIrange{0.6}{3.2}{\degree} RMS,
smaller than wheel odometry's, and explains at most one percent of the final
position error ($n=5$). The reported path length is wrong by
\SIrange{1}{249}{\percent}: the accelerometer integral drifts in speed, and
the position error follows a power law with exponent \SIrange{1.1}{1.75}{}
(median \num{1.5}, $R^2$ \SIrange{0.74}{0.94}{}), between the linear growth
of a velocity offset and the quadratic growth of an acceleration offset. The
gyroscope is not the weak sensor in this platform; the accelerometer
integration is. The growth is not concentrated in turns: on all five
recordings the position error grows faster on straights, where the speed
error accumulates over distance.

\FloatBarrier
\section{Insight 3: the EKF removes the heading error, and its covariance
         stops describing the error at the first turn}

The EKF reduces the heading error to \SIrange{0.3}{3.5}{\degree} RMS and the
position error by an order of magnitude, and its NEES leaves the consistent
band within one second of the first turn on every recording.
\Cref{tab:e1i-nees} and \cref{fig:nees} carry the timing;
\cref{tab:e1i-decomp} the decomposition.

\ins{nees_timing.tex}

\begin{figure}[htbp]
  \centering
  \includegraphics[width=\textwidth]{\insdir/fig_nees_light.png}
  \caption{3-DOF NEES of the EKF against time on each recording, for the
  measured per-recording noise model and for the VSLAM noise file, with the
  first ground-truth heading change (dotted), the \SI{97.5}{\percent} band
  (dashed) and the consistent value of 3 (solid).}
  \label{fig:nees}
\end{figure}

Before the first turn the filter is conservative: NEES sits below 3 on every
recording. Within \SI{0.6}{\second} of the first heading change it rises by
one to two orders of magnitude and never returns ($n=5$; departure at
\SIrange{9.0}{26.0}{\second} against first turns at
\SIrange{9.6}{25.6}{\second}). The covariance therefore describes the error
on straights and not the heading error that each turn injects: the process
model has no term for a turn-dependent heading error, so the state covariance
does not grow when the heading does. The two recordings on which the EKF is
least accurate, \repo{allsensor_003} and \repo{allsensor_004}, are the two
where its heading error shows a slow ramp between the turns of about
\SIrange{0.04}{0.07}{\degree\per\second} (\cref{fig:decomp}, bottom row),
which a gyroscope-bias state should absorb and here does not fully.

\begin{queried}[Consistency and calibration are different things]
The transferred noise model is less accurate and more consistent: its median
NEES is \SIrange{8}{39}{} against \SIrange{21}{91}{} for the measured model
($n=5$), because its larger process noise keeps the covariance closer to the
larger error. Neither model is consistent after the first turn. The raw
write-up's statement that the calibrated filter is ``overconfident by two
orders of magnitude'' is a whole-run average that mixes a consistent straight
with an inconsistent remainder.
\end{queried}

\FloatBarrier
\section{Insight 4: the transferred noise model turns the EKF back into
         wheel odometry at the turns}

With the VSLAM file's IMU constants the EKF accrues heading error at the turns
exactly as the raw wheel odometry does, and that is the whole of the damage
the noise transfer causes. The ``EKF @ VSLAM noise'' rows of
\cref{tab:e1i-decomp,tab:e1i-turning} carry the evidence.

On the transferred model the heading error grows at
\SIrange{431}{774}{\milli\degree\per\second} while turning and at or below
zero on straights, in five of five recordings; the position error grows at
\SIrange{59}{130}{\milli\metre\per\second} in the turns and at
\SIrange{-21}{8}{\milli\metre\per\second} between them. The signature is the
one of \cref{sec:insight1-wheel}: the loosened gyroscope noise density makes
the filter weight the wheel yaw rate over the gyroscope in the turns, and the
wheel yaw rate is the quantity that is wrong there. The raw write-up's
ablation found the gyroscope density to be the responsible term; this is the
mechanism by which it acts. Its path-length error of
\SIrange{9}{22}{\percent} is a secondary effect of the same rotation: the
along-track axis of the path frame moves with the true heading, so a heading
error puts part of the lateral error on it.

\FloatBarrier
\section{Insight 5: the visual estimators' reported error is largely
         vertical, and their one failure is a loop-closure event}

Scored in the plane, the visual estimators are better than the chapter's
three-dimensional ATE says, in four of fifteen runs by a factor of five or
more; their planar error does not grow at turns or with fewer features in any
consistent way; and the one gross failure is a single \SI{30}{\metre}
relocalisation at \SI{88}{\second} with the loop closer on, on a recording the
same estimator tracks to \SI{0.21}{\metre} with it off. \Cref{tab:e1i-visual}
and \cref{tab:e1i-orbfail} and \cref{fig:orb} carry the evidence.

\ins{visual.tex}

\begin{figure}[htbp]
  \centering
  \includegraphics[width=\textwidth]{\insdir/fig_orb_failure_light.png}
  \caption{ORB-SLAM3 with the loop closer on \repo{allsensor_003}. Left:
  whole-run aligned planar error against time, with the impossible-speed jumps
  of the estimate in red. Right: the aligned estimate against the ground truth
  at the IMU point on the warehouse plan, with the jump locations.}
  \label{fig:orb}
\end{figure}

The vertical component dominates where the three-dimensional ATE is large on
a run that tracked: VINS-Fusion on \repo{allsensor_000} is \SI{0.557}{\metre}
in three dimensions and \SI{0.115}{\metre} in the plane (vertical RMS
\SI{0.545}{\metre}); ORB-SLAM3 on the same recording \SI{0.507}{\metre}
against \SI{0.060}{\metre}; ORB-SLAM3 on \repo{allsensor_004}
\SI{1.171}{\metre} against \SI{0.219}{\metre}, and without the loop closer
\SI{0.756}{\metre} against \SI{0.135}{\metre}. The robot never leaves the
floor, so the vertical drift of a visual-inertial estimate is an error the
navigation task does not see, while the proprioceptive filter is planar by
construction and pays no such penalty. The chapter's estimator comparison
therefore charges the visual estimators for a component the baselines cannot
have.

\begin{queried}[Consequence for the Experiment 1 ranking]
In the plane the EKF is the most accurate estimator on one recording of five
(\repo{allsensor_002}) and the margins of the other four widen: on
\repo{allsensor_000} every visual configuration is below
\SI{0.12}{\metre} against the filter's \SI{0.19}{\metre}, on
\repo{allsensor_001} both estimators are below \SI{0.09}{\metre} against
\SI{0.31}{\metre}. Which ATE the chapter quotes for a floor robot is a
decision the methodology has to state; the planar value is the one that
matches the task.
\end{queried}

Turning and feature count do not organise the planar drift. VINS-Fusion
drifts faster while turning on two recordings of five and ORB-SLAM3 on three
of four with a usable regime; the Spearman correlation between matched
features and drift rate is between $-0.16$ and $0.26$ on every run. With one
run per estimator per recording neither is a pattern.

The failure is one event. ORB-SLAM3 with the loop closer on
\repo{allsensor_003} tracks for \SI{88}{\second} with no frame outside the
OK state and a planar error below \SI{4}{\metre} under the whole-run fit (the
fit itself is pulled by what follows), then takes jumps of \SI{30.0}{},
\SI{11.6}{} and \SI{2.5}{\metre} within \SI{1.5}{\second}. The net error
change inside the jump steps is \SI{+13.6}{\metre}; outside them it is
\SI{-2.9}{\metre}. The run without the loop closer on the same recording ends
at \SI{0.21}{\metre}. The event is therefore a loop closure or map merge that
accepted a wrong correspondence, late in the route, not tracking loss and not
drift; the raw write-up's flag on this value can be replaced by that
attribution.

\FloatBarrier
\section{What these data cannot answer}

\begin{itemize}
\item \textbf{Why the per-turn heading error of wheel odometry varies in
      size.} It is not the fitted yaw-rate scale alone (H1 of the raw
      write-up); slip at the steering limit, the Ackermann model's track width
      and the replay rate are candidates none of the logs isolate.
\item \textbf{Whether the EKF's residual heading ramp is gyroscope bias.} The
      bias state is logged in variance only, not in value, so the ramp on
      \repo{allsensor_003} and \repo{004} is attributed to it by elimination.
\item \textbf{Whether the loop-closure failure repeats.} One run with the loop
      closer on exists per recording; ORB-SLAM3 is non-deterministic.
\item \textbf{Anything about turning or features for the visual estimators
      with confidence.} Five recordings, one run each.
\end{itemize}

\section{Summary}

\begin{table}[htbp]
\centering\small
\caption{Questions, answers and the strength of the evidence.}
\label{tab:e1i-summary}
\begin{tabularx}{\textwidth}{lXXl}
\toprule
 & Answer & Evidence & Confidence \\
\midrule
Q1 wheel error & Heading error carried into position; distance within \SI{1}{\percent} & heading-only share 0.95--1.04 on 5/5 & high \\
Q2 turning & Wheel and transferred-noise EKF: heading made in turns, 5/5; IMU: no; visual: no pattern & \cref{tab:e1i-turning} & high / high / low \\
Q3 IMU divergence & Velocity integral, exponent $\approx 1.5$; heading explains $\le 1$ \% & \cref{tab:e1i-imu} & high \\
Q4 EKF consistency & Consistent until the first turn, inconsistent from it, 5/5 within \SI{1}{\second} & \cref{tab:e1i-nees} & high \\
Q5 noise transfer & Re-weights the wheel yaw rate in turns; heading error at turns like raw wheel odometry & ``EKF @ VSLAM noise'' rows & high \\
Q6 visual error & Vertical drift dominates 4/15 three-dimensional ATEs; planar error 0.04--0.7 m; one loop-closure event & \cref{tab:e1i-visual,tab:e1i-orbfail} & high (composition), medium (event) \\
\bottomrule
\end{tabularx}
\end{table}

\paragraph{Consequences for Chapter 4.} State which ATE a floor robot is
scored on and quote the planar value beside the three-dimensional one for the
visual estimators. Report the EKF's consistency as ``until the first turn''
rather than as a whole-run ANEES. Replace H1's falsification with the
mechanism: per-turn heading error of varying size. Attribute the ORB-SLAM3
\SI{6.4}{\metre} value to a single late loop-closure event rather than
flagging it as unexplained. Name the accelerometer integral, not the
gyroscope, as the reason IMU-only dead reckoning is unusable, since that is
what a VIO fusion is being asked to fix.

\end{document}
